\documentclass[11pt]{article}
\usepackage{amsmath,amssymb,amsthm}
\usepackage[margin=1in]{geometry}

\newtheorem{theorem}{Theorem}
\newtheorem{lemma}[theorem]{Lemma}
\newtheorem{corollary}[theorem]{Corollary}
\newtheorem{identity}[theorem]{Identity}
\theoremstyle{remark}
\newtheorem{remark}[theorem]{Remark}

\newcommand{\emd}{\operatorname{EMD}}
\newcommand{\qbin}[2]{\begin{bmatrix}#1\\#2\end{bmatrix}_q}

\title{The Rotation Orbit Tower and a Proof of $h_m\ge 0$ for $d=2$\\
\large Seed 4, Layer 2 --- Warnaar positivity project}
\author{Seed 4 (Layer 2, Round 2)}
\date{2026-07-04}

\begin{document}
\maketitle

\begin{abstract}
We prove two results toward the core bottleneck of the positivity
conjecture for $Q_{n,c}(q)$ ($r=3$, $d\not\equiv 0 \bmod 3$), namely the
conjecture $h_m := (q^{\ell};q^{\ell})_m\, g_m \ge 0$ with
$g_m = P_m/(q^3;q^3)_m$ and $\ell=\gcd(d,3)=1$.
First, an \emph{orbit-tower reduction} valid for all $d$ with $3\nmid d$:
$h_m$-positivity is equivalent to coefficientwise nonnegativity of an
explicit tower $H_m = U(q^m)H_{m-1}$, $H_0=\mathbf 1$, where $U(x)$ is a
fixed $N\times N$ polynomial matrix, $N=\tfrac{(d+1)(d+2)}{6}$, indexed by
the $\mathbb Z/3$-rotation orbits of profiles. The exact divisibility making
$U$ polynomial follows from a two-line congruence for the cyclic earth-mover
distance. Second, for $d=2$ we solve the tower in closed form: $h_m$ equals
a finite Rogers--Ramanujan (MacMahon--Schur) polynomial, hence is manifestly
nonnegative. The proof uses only the $q$-Pascal and absorption identities.
\end{abstract}

\section{Setup}

Fix $r=3$ and $d\ge 1$ with $3\nmid d$. A \emph{profile} is a composition
$c=(c_0,c_1,c_2)$ of $d$ into three nonnegative parts; there are
$\binom{d+2}{2}$ of them. The cyclic (clockwise) earth-mover distance is
\[
\emd(c,c') \;=\; 3\max(0,\alpha,\beta)-\alpha-\beta,
\qquad \alpha = c'_1-c_1,\quad \beta = c_0-c'_0 .
\]
Define $P_m(c)$ as the generating function of weighted profile paths,
\[
P_0(c)=1,\qquad
P_m(c)\;=\;\sum_{c'} q^{\,m\,\emd(c,c')}\,P_{m-1}(c'),
\]
so that $P_m(c)=\sum_{\gamma} q^{\sum_{k=1}^m k\,\emd(c^{k},c^{k-1})}$ over
paths $\gamma=(c^0,\dots,c^m=c)$. The Layer-1 results (Adjugate Monomial
Theorem and its consequences, all verified) express the core quantity as
\[
h_m(c)\;=\;\frac{(q;q)_m}{(q^3;q^3)_m}\,P_m(c)
\;=\;\frac{P_m(c)}{\prod_{k=1}^m (1+q^k+q^{2k})},
\]
using $(1-q^{3k})/(1-q^k)=1+q^k+q^{2k}$. The bottleneck conjecture is that
$h_m(c)$ is a polynomial with nonnegative coefficients.

Let $\rho(c)=(c_2,c_0,c_1)$ denote cyclic rotation of profiles. Directly
from the formula, $\emd(\rho a,\rho b)=\emd(a,b)$; hence by induction
$P_m(\rho c)=P_m(c)$ and $h_m$ is constant on $\rho$-orbits. Since
$3\nmid d$, no profile is $\rho$-fixed, so all orbits have size $3$ and
there are $N=\binom{d+2}{2}/3=\frac{(d+1)(d+2)}{6}$ of them.

\section{The orbit tower}

\begin{lemma}[mod-3 congruence for EMD]\label{lem:mod3}
For all profiles $a,b$ of $d$,
\[
\emd(a,b)\;\equiv\; b_0-a_0+a_1-b_1 \pmod 3,
\qquad
\emd(a,\rho b)-\emd(a,b)\;\equiv\; d \pmod 3 .
\]
\end{lemma}

\begin{proof}
Since $\emd = 3\max(0,\alpha,\beta)-\alpha-\beta \equiv -\alpha-\beta
= (a_1-b_1)+(b_0-a_0) \pmod 3$, the first congruence holds. With
$\rho b=(b_2,b_0,b_1)$,
\[
\bigl[(\rho b)_0-a_0+a_1-(\rho b)_1\bigr]-\bigl[b_0-a_0+a_1-b_1\bigr]
= (b_2-b_0)-(b_0-b_1)= b_0+b_1+b_2-3b_0 \equiv d \pmod 3. \qedhere
\]
\end{proof}

\begin{theorem}[Orbit-tower reduction, all $3\nmid d$]\label{thm:tower}
Choose orbit representatives $a_1,\dots,a_N$, and for each pair $(i,j)$ set
$E_t(i,j)=\emd(a_i,\rho^t a_j)$, $t=0,1,2$. Then:
\begin{enumerate}
\item The exponents $E_0,E_1,E_2$ are pairwise distinct modulo $3$; hence
$x^{E_0}+x^{E_1}+x^{E_2}$ is divisible by $1+x+x^2$ in $\mathbb Z[x]$, and
\[
U_{ij}(x)\;:=\;\frac{x^{E_0(i,j)}+x^{E_1(i,j)}+x^{E_2(i,j)}}{1+x+x^2}
\;\in\;\mathbb Z[x].
\]
\item The vector $H_m=(h_m(a_1),\dots,h_m(a_N))$ satisfies
\[
H_m \;=\; U(q^m)\,H_{m-1},\qquad H_0=(1,\dots,1)^{\mathsf T}.
\]
\end{enumerate}
Consequently $h_m\ge 0$ (all $c$, all $m$) if and only if the tower
$H_m$ is coefficientwise nonnegative for all $m$.
\end{theorem}

\begin{proof}
(1) By Lemma~\ref{lem:mod3}, $E_t\equiv E_0+td \pmod 3$; since $3\nmid d$
these are the three residues in some order. If $\omega$ is a primitive cube
root of unity, $\omega^{E_0}+\omega^{E_1}+\omega^{E_2}=1+\omega+\omega^2=0$,
so both primitive cube roots of unity are roots of
$x^{E_0}+x^{E_1}+x^{E_2}$, which is therefore divisible by $1+x+x^2$.

(2) Group the $P$-recursion by orbits, using $\rho$-invariance of
$P_{m-1}$:
\[
P_m(a_i)=\sum_{j=1}^N\Bigl(\sum_{t=0}^2 q^{m E_t(i,j)}\Bigr)P_{m-1}(a_j).
\]
Divide both sides by $\prod_{k\le m}(1+q^k+q^{2k})$; the factor
$1+q^m+q^{2m}$ absorbs into the bracket by part (1) with $x=q^m$, giving
the stated system for $H_m$. The base case is $P_0=1$.
\end{proof}

\begin{remark}
Lemma~\ref{lem:mod3} turns the equidistribution phenomenon observed
computationally in earlier layers (Seed 7) into a proof. Theorem
\ref{thm:tower} replaces the infinite family of statements
``$(q;q)_m P_m/(q^3;q^3)_m\ge 0$'' by nonnegativity of a tower generated
by one \emph{fixed, explicitly computable} polynomial matrix.
\end{remark}

\section{Solution of the tower for $d=2$}

For $d=2$ there are $N=2$ orbits, with representatives $a_1=(2,0,0)$ and
$a_2=(1,1,0)$. Direct computation of the twelve values
$\emd(a_i,\rho^t a_j)$ gives the exponent triples
\[
\{E_t(1,1)\}=\{0,2,4\},\quad
\{E_t(1,2)\}=\{E_t(2,1)\}=\{1,2,3\},\quad
\{E_t(2,2)\}=\{0,1,2\},
\]
so, since $(1+x^2+x^4)=(1+x+x^2)(1-x+x^2)$,
\[
U(x)=\begin{pmatrix}1-x+x^2 & x\\ x & 1\end{pmatrix}.
\]

Recall the Gaussian binomial $\qbin{m}{j}$ and define the finite
Rogers--Ramanujan (MacMahon--Schur) polynomials
\[
A_m=\sum_{j\ge 0} q^{j^2+j}\qbin{m}{j},\qquad
B_m=\sum_{j\ge 0} q^{j^2}\qbin{m}{j},\qquad
C_m=\sum_{j\ge 0} q^{j^2+2j}\qbin{m}{j}.
\]

\begin{theorem}[$d=2$ solved]\label{thm:d2}
For $d=2$ and all $m\ge 0$,
\[
h_m(c)=A_m \ \ \text{if } c\in\{(2,0,0),(0,2,0),(0,0,2)\},
\qquad
h_m(c)=B_m \ \ \text{if } c\in\{(1,1,0),(0,1,1),(1,0,1)\}.
\]
In particular $h_m(c)$ has nonnegative coefficients for all $c$ and $m$.
\end{theorem}

The proof rests on three elementary identities. We use the $q$-Pascal rule
$\qbin{m}{j}=\qbin{m-1}{j}+q^{m-j}\qbin{m-1}{j-1}$ and the absorption
identities $(1-q^j)\qbin{n}{j}=(1-q^n)\qbin{n-1}{j-1}$ and
$(1-q^{n-j})\qbin{n}{j}=(1-q^n)\qbin{n-1}{j}$.

\begin{identity}\label{id:B}
$B_m-B_{m-1}=q^m A_{m-1}$.
\end{identity}
\begin{proof}
By $q$-Pascal, $B_m-B_{m-1}=\sum_j q^{j^2}\,q^{m-j}\qbin{m-1}{j-1}$.
Substituting $j=i+1$: $\;q^m\sum_i q^{i^2+i}\qbin{m-1}{i}=q^mA_{m-1}$.
\end{proof}

\begin{identity}\label{id:A}
$A_m-A_{m-1}=q^{m+1} C_{m-1}$.
\end{identity}
\begin{proof}
By $q$-Pascal, $A_m-A_{m-1}=\sum_j q^{j^2+j}\,q^{m-j}\qbin{m-1}{j-1}
= q^m\sum_i q^{(i+1)^2}\qbin{m-1}{i}=q^{m+1}C_{m-1}$.
\end{proof}

\begin{identity}[key lemma]\label{id:C}
$q\,C_n = B_n-(1-q^{n+1})A_n$ for all $n\ge 0$.
\end{identity}
\begin{proof}
Compute
\[
B_n-A_n+q^{n+1}A_n-qC_n
=\sum_j \qbin{n}{j}q^{j^2}\Bigl((1-q^j)-q^{2j+1}(1-q^{n-j})\Bigr).
\]
Apply the first absorption identity to the term $(1-q^j)\qbin{n}{j}$ and
the second to $(1-q^{n-j})\qbin{n}{j}$:
\[
=(1-q^n)\Bigl(\sum_j q^{j^2}\qbin{n-1}{j-1}
-\sum_j q^{j^2+2j+1}\qbin{n-1}{j}\Bigr)
=(1-q^n)\Bigl(\sum_i q^{i^2+2i+1}\qbin{n-1}{i}
-\sum_j q^{j^2+2j+1}\qbin{n-1}{j}\Bigr)=0,
\]
where the first sum was reindexed by $j=i+1$.
(For $n=0$ all three of $A_0,B_0,C_0$ equal $1$ and the identity reads
$q=1-(1-q)$.)
\end{proof}

\begin{proof}[Proof of Theorem~\ref{thm:d2}]
By Theorem~\ref{thm:tower} it suffices to show that
$(A_m,B_m)$ satisfies the $d=2$ tower, i.e.\ with $x=q^m$,
\[
A_m=(1-x+x^2)A_{m-1}+xB_{m-1},\qquad
B_m=xA_{m-1}+B_{m-1},
\]
with base case $A_0=B_0=1=H_0$. The $B$-equation is
Identity~\ref{id:B}. For the $A$-equation,
\[
(1-q^m+q^{2m})A_{m-1}+q^mB_{m-1}
=A_{m-1}+q^m\bigl(B_{m-1}-(1-q^m)A_{m-1}\bigr)
=A_{m-1}+q^m\cdot qC_{m-1}=A_m,
\]
using Identity~\ref{id:C} with $n=m-1$ and then Identity~\ref{id:A}.
Since the tower has a unique solution given $H_0$, we conclude
$H_m=(A_m,B_m)$, i.e.\ $h_m(a_1)=A_m$, $h_m(a_2)=B_m$, and by
$\rho$-invariance the same holds on the full orbits. Nonnegativity is
manifest from the defining sums.
\end{proof}

\paragraph{Example ($m=2$).}
$A_2=1+q^2+q^3+q^6$, $B_2=1+q+q^2+q^4$. Check the tower with $x=q^2$:
$B_2=q^2A_1+B_1=q^2(1+q^2)+(1+q)$ \checkmark;
$A_2=(1-q^2+q^4)(1+q^2)+q^2(1+q)=(1+q^6)+(q^2+q^3)$ \checkmark.
Also $\prod_{k\le 2}(1+q^k+q^{2k})\cdot A_2 = P_2((2,0,0))$, verified
against the direct path enumeration.

\section{Status for $d=4$ and outlook}

For $d=4$ the tower is $5\times 5$. Computationally, three of the five
orbit values admit bounded fermionic forms of Warnaar's Conjecture-2 shape
\[
H_m=\sum_{n,j} q^{n^2-nj+j^2+an+bj}\qbin{m}{n}\begin{bmatrix}2n\\
j\end{bmatrix}_q,
\]
namely the orbits of $(1,1,2)$, $(0,0,4)$, and $(0,3,1)$ (verified for
$m\le 5$). The two remaining orbits, of $(0,1,3)$ and $(0,2,2)$, provably
do not fit a single sum of this shape, and an exhaustive search over pairs
of such sums also failed; these profiles lie outside the family covered by
Warnaar's Conjecture~2 and finding their fermionic forms is open. The
$d=2$ proof provides the template: once closed forms are found for all
five orbits, the tower identities should again reduce to $q$-Pascal
manipulations.

\end{document}
